\chapter{Introduction}
This document outlines the concept for a "Gift List" application, designed to streamline the process of gift-giving for various celebrations and events. The core idea is to provide a platform where a celebrant (e.g., a birthday person, a couple getting married) can curate and share a personalized wish list of gifts they would like to receive.

\myskip

The application serves users who can switch between two primary roles within the platform:
\begin{itemize}
    \item As a Celebrant \\
        A user creates and manages their own gift lists. They can add, remove, and describe desired items. Critically, as the creator of a list, the user will not have visibility into the real-time status of their gift list regarding which items have been claimed. This maintains the element of surprise for them.
    \item As a Guest \\
        A user receives a link to another user's gift list. To view the items or interact with the list, they must be registered and logged into the platform. Once authenticated, they can browse the available items and "claim" a gift they intend to purchase. This mechanism prevents duplicate gifts and helps guests coordinate their purchases efficiently.
\end{itemize}

% LOGIC: All users must register to access the platform's features, allowing each user to act both as a list creator (celebrant) and a gift giver (guest). This enables a unified dashboard experience.
The primary goal of this application is to simplify gift coordination, reduce redundant gifts, and enhance the gift-giving experience by ensuring celebrants receive desired items while preserving the element of surprise, all within a unified platform.

\begin{figure}[htbp]
    \centering
    \includegraphics[width=.4\textwidth, trim=0 150 0 0, clip]{images/1-introduction/use-case.pdf}
    \caption{Use case diagram of Gift List requirements.}
\end{figure}

%%%%%%%%%%%%%%%%%%%%%%%%%%%%%%%%%%%%%%%%%%%%%%%%%%%%%%%%
\section{Requirements}
In this Section we will define the requirements for the application. The application must meet the following functional requirements:
\begin{itemize}
    \item \textbf{User Management and Dashboard}
    \begin{itemize}
        \item User Authentication \\
            The application must require all users to register and log-in to access any features. Whether creating a new list or responding to an invitation, an authenticated account is mandatory.
        \item Unified Dashboard \\
            Users must have access to a centralized dashboard divided into two distinct sections: lists they have created (acting as a Celebrant) and lists they have joined (acting as a Guest).
        \item List Management \\
            Users must be able to create and delete their own gift lists from the dashboard.
        \item Item Management \\
            For lists they own, users must be able to add, edit, and remove items. For each item, they can provide a name, description, an optional external URL, and set a preference level.
    \end{itemize}

    \item \textbf{Guest Interaction}
    \begin{itemize}
        \item List Access \\
            Users must be able to view a specific gift list using a unique URL invitation. Accessing this URL will either prompt login/registration or immediately link the list to their account, making it visible in their "joined lists" dashboard section.
        \item Item Claiming \\
            Users acting as guests must be able to mark an item as "claimed" or "purchased". This action must be visible to other guests to prevent duplicate purchases but must remain hidden from the list owner to preserve the surprise.
        \item Item Unclaiming \\
            Users must be able to revert a claim, making the item available again. This action is only permitted for the user who originally claimed the item.
        \item Personal Claim Visibility \\
            Users viewing a joined list should clearly distinguish which items they have personally claimed versus those claimed by others.
    \end{itemize}

    \item \textbf{System Integrity}
    \begin{itemize}
        \item Claimed Item Removal \\
            If a Celebrant removes an item that has already been claimed by a Guest, the system must allow the removal but must immediately notify the Guest who claimed it. This prevents the Guest from purchasing an item that is no longer desired, while maintaining the Guest's anonymity.
    \end{itemize}
\end{itemize}

%%%%%%%%%%%%%%%%%%%%%%%%%%%%%%%%%%%%%%%%%%%%%%%%%%%%%%%%
\section{Internationalization (i18n)}
The application is designed to be accessible to a global audience, with initial support for two languages:
\begin{itemize}
    \item \textbf{English}: The default language for the interface and system communications.
    \item \textbf{Italian}: Full support for the Italian language.
\end{itemize}
The application will detect the user's browser language setting to provide the most appropriate initial experience, while also allowing users to manually switch between supported languages via a language selector in the UI.

%%%%%%%%%%%%%%%%%%%%%%%%%%%%%%%%%%%%%%%%%%%%%%%%%%%%%%%%
\section{Accessibility}
The "Gift List" application will follow a phased release strategy regarding platform availability:
\begin{itemize}
    \item Initial Phase: Web Application \\
        The application will initially be accessible exclusively as a web application through modern web browsers. Official support will be focused on the most widely used browsers (e.g., Google Chrome, Mozilla Firefox, Apple Safari, and Microsoft Edge). Limiting official support to these "evergreen" browsers ensures a consistent user experience, optimal performance, and robust security by leveraging the latest web standards and security protocols.
    \item Future Phase: Mobile Applications \\
        Following the stabilization of the web platform, dedicated native or hybrid applications are planned for distribution through the Apple App Store and Google Play Store to provide a more integrated mobile experience.
\end{itemize}

%%%%%%%%%%%%%%%%%%%%%%%%%%%%%%%%%%%%%%%%%%%%%%%%%%%%%%%%
\section{Licensing}
The "Gift List" project is released as an open-source application under the \textbf{MIT License}. This permissive license allows for the free use, modification, and distribution of the software, provided that the original copyright notice and permission notice are included in all copies or substantial portions of the software. This choice supports the project's philosophy of community access and collaborative improvement.